\documentclass[10pt,a4paper]{article}
%\usepackage[utf8x]{inputenc}

%\usepackage[latin1]{inputenc}
%\usepackage[spanish]{babel}

\usepackage{amssymb,latexsym,amsthm,amsmath}
\usepackage{verbatim}
\usepackage{graphicx,color,wrapfig}
\usepackage{psfrag}
\usepackage[all]{xy}
\usepackage{tikz}
%\includepackage{mathabx}
\usepackage{hyperref}
\synctex=1
\usepackage{srcltx}
\hypersetup{colorlinks=true,linkcolor=blue}

\sloppy
\usepackage{multicol}
\setlength{\topmargin}{-0.99in}
\setlength{\textheight}{9.95in}
\setlength{\oddsidemargin}{-0.3in}
\setlength{\evensidemargin}{-0.3in}
\setlength{\textwidth}{6.8in}

\def\R{\mathbb R}
\def\Q{\mathbb Q}
\def\N{\mathbb N}

\DeclareMathOperator{\dom}{dom}
\DeclareMathOperator{\ran}{ran}

\definecolor{light-gray}{gray}{0.85}


\usepackage{type1cm}
\usepackage{eso-pic}
\usepackage{calc}
\usepackage{forloop}

\newcounter{ber}
\newcounter{num}

\newcommand{\lef}{\langle}
\newcommand{\rig}{\rangle}

\makeatletter

\makeatother


\begin{document}

\begin{center}
\begin{figure}[ht!]
\begin{minipage}[c]{0.2\linewidth}
\includegraphics[scale=0.22]{EscudoUNI.eps}
\end{minipage}
\hfill{}
\begin{minipage}{0,8\columnwidth}
\centering\flushleft {\large\bf Introducci\'on a la Programaci\'on\\
Escuela Profesional de Ciencia de la Computaci\'on\\
Facultad de Ciencias\\
Universidad Nacional de Ingenier\'ia\\}
\end{minipage}
\end{figure}
\end{center}
\vspace*{-1.2cm}
\noindent

\begin{center}
\textbf{\Large Laboratorio dirigido 5}
\end{center}
\textbf{Tema: Paso de par\'ametros, alcance de variables y recursividad \hfill Ciclo: 2015-I}
\noindent
\rule{\textwidth}{1pt}
\newline
%\vspace*{0.5cm}



Empleando en cada ejercicio la t\'ecnica de recursividad, implemente un programa en el que pida ingresar
\begin{enumerate}
\item Un entero positivo $n$ y muestre el valor de la siguiente suma: $$ 1^1 - 2^2 + 3^3 - 4^4 + \ldots n^n. $$
\item Un entero positivo $n$ y muestre el $n$-\'esimo t\'ermino de la sucesi\'on de Fibonacci.
\item Un margen de error ($\epsilon > 0$) y muestre el valor de la primera suma parcial que con un margen de error de $\epsilon$ se aproxime al valor de la siguiente serie: 
$$1 - \frac{1}{2} + \frac{1}{3} - \frac{1}{4} + \cdots = \ln(2).$$
\item Un margen de error ($\epsilon > 0$) y muestre el valor de la primera suma parcial que con un margen de error de $\epsilon$ se aproxime al valor de la siguiente serie: 
$$\frac{1}{1^2} + \frac{1}{2^2} + \frac{1}{3^2} + \cdots = \frac{{\pi}^2}{6}.$$
\item Una cadena de caracteres $s$ y devuelva todas las combinaciones posibles que se pueden formar a partir de los caracteres de $s$, e.g. entrada: ``sol''; salida: ``sol'', ``slo'', ``osl'', ``ols'', ``lso'', ``los''.
\item Un n\'umero positivo $b$ y un entero $n$. Luego, muestre el valor de $b^n$.
\item Dos enteros positivos $a$ y $b$. Luego, muestre el MCD de $a$ y $b$. (Sugerencia: puede usar el Algoritmo de Euclides para calcular el MCD.)
\item Un entero positivo $n$ y devuelva el n\'umero que resulta de invertir el orden de las cifras de $n$. (Sugerencia: puede emplear una variable global para almacenar el n\'umero a mostrar.)
\item Un entero positivo $n$ y un entero positivo $k$ menor o igual a $n$. Luego, calcule el n\'umero de combinaciones de $k$ elementos de $n$ elementos dados. (Sugerencia: puede utilizar la identidad abajo mostrada.) $$\binom nk = \binom{n-1}{k-1} + \binom{n-1}k.$$   

\end{enumerate}

\end{document}
